\section{Lift selections of low degeneracy}\label{sec:lifts}

This section is pure graph theory. It produces, on a finite regular graph, the structure that \cref{sec:peeling} turns
into private checks on the tree.

\subsection{Lift selections}\label{sec:lifts-def}

Let $X$ be a finite simple graph. Its \emph{lifted vertices} are the pairs $(v,i)$ with $v\in X^0$ and $i\in\{0,1\}$,
and the \emph{swap} is $\sigma(v,i)=(v,1-i)$. An edge $e=\{x,y\}$ of $X$ has four \emph{lifts}
$\{(x,i),(y,j)\}$, exchanged in pairs by~$\sigma$.

\begin{definition}[Lift selection]\label{def:selection}
\formal{AboveHalfRate/LiftSelection.lean}{AboveHalf.IsLift,AboveHalf.isLift_iff,AboveHalf.liftSwap,AboveHalf.IsLift.map_liftSwap,AboveHalf.IsEarlierEnd,AboveHalf.IsDegenerateFor,AboveHalf.liftPos_injective}
A \emph{lift selection} $A$ of $X$ consists of one lift $\tilde a(e)$ of each edge $e$ of $X$. Its \emph{complement} is
$B=\sigma(A)$, the set of lifts $\sigma(\tilde a(e))$. A lift selection is \emph{$t$-degenerate} if there is a linear
order of the lifted vertices in which every lifted vertex is the earlier endpoint of at most $t$ lifts in~$A$.
\end{definition}

In Lean, $t$-degeneracy is stated for the order given by real positions of the lifted vertices; it is the notion of
\cref{def:selection} when the positions are distinct, as the positions we use are (\lean{AboveHalf.liftPos_injective}).

The lifts in $A\cup B$ form a graph $\Ftil X$ on the lifted vertices, and $(v,i)\mapsto v$ is a two-sheeted covering
$\Ftil X\to X$: for each edge $e$ at $x$, the lifted vertex $(x,i)$ lies in exactly one of $\tilde a(e)$ and
$\sigma(\tilde a(e))$. The swap is a deck transformation exchanging $A$ and $B$, so $B$ is $t$-degenerate whenever $A$
is, for the order transported by~$\sigma$. A lift selection is the same as a voltage assignment in $\Z/2$ together with
a choice, for each edge, of which of its two lifts in the double cover is called~$A$.

\begin{theorem}[Theorem D]\label{thm:D}
\partformal{AboveHalfRate/LiftSelection.lean}{AboveHalf.exists_liftSelection_of_fluid,AboveHalf.liftPos_injective}{the probabilistic existence of labels with small loads (\cref{lem:concentration}),
which uses Bernstein's inequality.}
Let $X$ be a finite simple graph on $n$ vertices, regular of degree $\delta\ge1$. Put $\tau=2-\sqrt3$ and
$\nu=\tau(1-\tau)\delta$, and let $\lambda>0$ satisfy
\begin{equation}\label{eq:Dcond}
  2n\exp\Bigl(-\frac{\lambda^2}{2(\nu+\lambda/3)}\Bigr)<1.
\end{equation}
Then $X$ has a lift selection that is $t$-degenerate for $t=\lceil\tau\delta+\lambda\rceil$. The order can be taken as
follows, for injective labels $\xi\colon X^0\to(0,1)$: the lifted vertex $(v,0)$ has position $\xi_v$, the lifted vertex
$(v,1)$ has position $2-\xi_v$, and lifted vertices are ordered by position.
\end{theorem}
\begin{outline}
The formal part is \lean{AboveHalf.exists_liftSelection_of_fluid}. From injective labels in $(0,1)$ whose fluid loads
(\cref{sec:lifts-fluid}) are at most an integer $t$, it produces a lift selection $A$ such that $A$ is $t$-degenerate
for the order by position and $\sigma(A)$ for its transport by~$\sigma$. It rounds the fluid rule by Hall's theorem
(\cref{lem:rounding}) and realizes the chosen options by lifts (\cref{lem:options}); \lean{AboveHalf.liftPos_injective}
says that the positions are distinct. The existence of labels with loads at most $\lceil\tau\delta+\lambda\rceil$ is the
probabilistic step, proved in \cref{sec:lifts-proof}.
\end{outline}

Since $\lambda=O(\sqrt{\delta\log n}+\log n)$ satisfies \eqref{eq:Dcond}, the degeneracy is
$(2-\sqrt3)\delta+O(\sqrt{\delta\log n})$ when $\log n=O(\delta)$. The constant $2-\sqrt3\approx0.268$ is what the
construction of \cref{sec:theoremA} pays: local codes of redundancy $t$ carry the separated structure, and the rate
obtained is about $(1-t/\delta)^2$. A naive selection does much worse; \cref{rem:naive} explains why some cleverness is
needed.

\subsection{Options and rounding}\label{sec:lifts-options}

Fix injective labels $\xi\colon X^0\to(0,1)$ and order the lifted vertices by position as in \cref{thm:D}. All first
copies $(v,0)$ come before all second copies $(v,1)$; the first copies are in increasing order of label, the second
copies in decreasing order.

\begin{lemma}[Options]\label{lem:options}
\formal{AboveHalfRate/LiftSelection.lean}{AboveHalf.isEarlierEnd_lifts,AboveHalf.isOption_mk_iff,AboveHalf.isOption_iff_exists_lift,AboveHalf.liftOf_options,AboveHalf.isEarlierEnd_liftOf,AboveHalf.isDegenerateFor_image_liftOf}
Let $e=\{x,y\}$ be an edge with $\xi_x<\xi_y$. The earlier endpoints of its four lifts are
\begin{gather*}
  \{(x,0),(y,0)\}\mapsto(x,0),\qquad \{(x,0),(y,1)\}\mapsto(x,0),\\
  \{(x,1),(y,0)\}\mapsto(y,0),\qquad \{(x,1),(y,1)\}\mapsto(y,1).
\end{gather*}
Call $O(e)=\{(x,0),(y,0),(y,1)\}$ the \emph{options} of~$e$. For every map $o$ that assigns to each edge one of its
options, the lift selection
\[
  \tilde a(e)=\{(x,0),(y,1)\},\ \{(x,1),(y,0)\},\ \{(x,1),(y,1)\}\quad\text{for}\quad o(e)=(x,0),\ (y,0),\ (y,1)
\]
has earlier endpoint $o(e)$ on every edge. Hence it is $t$-degenerate if $o$ chooses every lifted vertex for at most
$t$ edges.
\end{lemma}
\begin{outline}
The positions satisfy $\xi_x<\xi_y<1<2-\xi_y<2-\xi_x$, which gives the four earlier endpoints. The construction realizes
each option by a lift whose earlier endpoint is that option.
\end{outline}

A \emph{fractional assignment} gives each edge $e$ nonnegative weights $w(e,z)$ on its options $z\in O(e)$, with sum~$1$.
The \emph{load} of a lifted vertex $z$ is $\sum_ew(e,z)$, the sum over the edges having $z$ as an option.

\begin{lemma}[Rounding]\label{lem:rounding}
\formal{AboveHalfRate/LiftSelection.lean}{AboveHalf.exists_assignment_of_fractional}
If a fractional assignment has all loads at most an integer $T$, then some assignment of options has every lifted
vertex chosen by at most $T$ edges.
\end{lemma}
\begin{outline}
Make $T$ copies of each lifted vertex and join each edge to the copies of its options. For a set $S$ of edges, the
total weight $\abs S$ is spread over the options of $S$, each receiving at most $T$, so the copies of those options
number at least $\abs S$. By Hall's theorem~\cite{Hall1935} the edges can be matched to distinct copies.
\end{outline}

\subsection{The fluid rule}\label{sec:lifts-fluid}

Put $\tau=2-\sqrt3$, the smaller root of $\tau^2-4\tau+1=0$; equivalently $(1-\tau)^2=2\tau$. For an edge $\{x,y\}$ with
$\xi_x<\xi_y$, write $\eta=\xi_y$ and send
\[
  \min\Bigl(1,\frac\tau\eta\Bigr)\ \text{to}\ (y,1),\qquad
  \Bigl(1-\frac\tau\eta\Bigr)_+\frac{\eta-\tau}{2\eta}\ \text{to}\ (y,0),\qquad
  \Bigl(1-\frac\tau\eta\Bigr)_+\frac{\eta+\tau}{2\eta}\ \text{to}\ (x,0).
\]
The weights are nonnegative and sum to~$1$. They depend only on the larger label. For a vertex $z$ with label $u$ and a
neighbour with label $s\ne u$, the weight that the edge between them sends to $(z,1)$ and to $(z,0)$ is
\[
  f_1(u,s)=\begin{cases}\min(1,\tau/u)&s<u,\\0&s>u,\end{cases}\qquad
  f_0(u,s)=\begin{cases}(1-\tau/u)_+\,\dfrac{u-\tau}{2u}&s<u,\\[1ex](1-\tau/s)_+\,\dfrac{s+\tau}{2s}&s>u,\end{cases}
\]
so the load of $(z,i)$ is $\sum_{w\sim z}f_i(u,\xi_w)$.

\begin{lemma}[Fluid loads]\label{lem:fluid}
\formal{AboveHalfRate/Fluid.lean}{AboveHalf.Fluid.charges_sum,AboveHalf.Fluid.f1_nonneg,AboveHalf.Fluid.f1_le_one,AboveHalf.Fluid.f0_nonneg,AboveHalf.Fluid.f0_le_half,AboveHalf.Fluid.integral_f1,AboveHalf.Fluid.integral_f0,AboveHalf.Fluid.one_sub_tau_sq,AboveHalf.Fluid.integral_f1_sq_sub_le,AboveHalf.Fluid.integral_f0_sq_le,AboveHalf.Fluid.integral_f0_sq_sub_le}
For every $u\in(0,1]$:
\begin{enumerate}[label=(\alph*)]
\item $0\le f_1(u,s)\le1$ and $0\le f_0(u,s)\le\frac12$ for $s\in(0,1]$;
\item $\int_0^1f_1(u,s)\,ds=\min(u,\tau)\le\tau$ and $\int_0^1f_0(u,s)\,ds=\frac{(1-\tau)^2}2=\tau$;
\item $\int_0^1f_1(u,s)^2\,ds-\bigl(\int_0^1f_1(u,s)\,ds\bigr)^2\le\tau(1-\tau)$ and
$\int_0^1f_0(u,s)^2\,ds\le\tau/2\le\tau(1-\tau)$.
\end{enumerate}
\end{lemma}
\begin{outline}
(a) is direct. (b) For $f_1$ the integral is $u\min(1,\tau/u)$. For $f_0$ it is
$(u-\tau)_+^2/(2u)+\frac12\int_{\max(u,\tau)}^1(1-\tau^2/s^2)\,ds$. When $u\le\tau$ this is
$\frac12(1-\tau)-\frac12\tau(1-\tau)=\frac12(1-\tau)^2$. When $u>\tau$ it is
$\bigl((u-\tau)^2+u(1-u)-\tau^2(1-u)\bigr)/(2u)=\frac12(1-\tau)^2$. Finally $(1-\tau)^2=2\tau$. (c) For $u\le\tau$ the
first variance is $u(1-u)\le\tau(1-\tau)$; for $u>\tau$ it is $\tau^2(1-u)/u\le\tau(1-\tau)$. The second bound follows
from $f_0\le\frac12$ and~(b).
\end{outline}

So each lifted vertex receives load $\tau$ per unit degree in expectation, or less, whatever its label: the rule is
balanced exactly at $\tau=2-\sqrt3$.

\subsection{Concentration and the proof of Theorem D}\label{sec:lifts-proof}

We use Bernstein's inequality~\cite[Section~2.8]{BLM2013}: if $Y_1,\dots,Y_\delta$ are independent random variables
with $Y_i-\mathbb EY_i\le M$ and $\sum_i\operatorname{Var}Y_i\le\nu$, then
$\Pr\bigl(\sum_i(Y_i-\mathbb EY_i)\ge\lambda\bigr)\le\exp\bigl(-\lambda^2/(2(\nu+M\lambda/3))\bigr)$ for $\lambda>0$.

\begin{lemma}[Concentration]\label{lem:concentration}
\notformalized
Let the labels $\xi_v$ be independent and uniform on $(0,1)$. For every lifted vertex $(z,i)$ and every $\lambda>0$,
the load of $(z,i)$ under the fluid rule is at least $\tau\delta+\lambda$ with probability at most
$\exp\bigl(-\lambda^2/(2(\nu+\lambda/3))\bigr)$, where $\nu=\tau(1-\tau)\delta$.
\end{lemma}
\begin{informalproof}
Condition on $\xi_z=u$. Since $X$ is simple, the $\delta$ neighbours $w$ of $z$ are distinct vertices other than $z$, so
their labels are independent and uniform on $(0,1)$ given $\xi_z$. The load is $\sum_{w\sim z}f_i(u,\xi_w)$, a sum of
$\delta$ independent terms with values in $[0,1]$. By \cref{lem:fluid}(b) its conditional mean is at most $\tau\delta$,
and by \cref{lem:fluid}(c) the conditional variances sum to at most $\nu$. Bernstein's inequality with $M=1$ bounds the
conditional probability that the load exceeds its mean by $\lambda$, hence the conditional probability that it is at
least $\tau\delta+\lambda$, by $\exp(-\lambda^2/(2(\nu+\lambda/3)))$. Integrating over $u$ gives the claim.
\end{informalproof}

\begin{proof}[Proof of \cref{thm:D}]
Choose the labels independently and uniformly on $(0,1)$. They are injective with probability $1$, and by
\cref{lem:concentration} and a union bound over the $2n$ lifted vertices, all loads are below $\tau\delta+\lambda$ with
probability at least $1-2n\exp(-\lambda^2/(2(\nu+\lambda/3)))>0$. Fix injective labels with this property. The fluid rule is a
fractional assignment with all loads at most the integer $t=\lceil\tau\delta+\lambda\rceil$. By \cref{lem:rounding} some
assignment of options has every lifted vertex chosen by at most $t$ edges, and by \cref{lem:options} the corresponding
lift selection is $t$-degenerate for the order by position.
\end{proof}

Only the probabilistic step is not formalized. Given injective labels in $(0,1)$ whose fluid loads are at most $t$,
\lean{AboveHalf.exists_liftSelection_of_fluid} produces a lift selection $A$ such that $A$ and $\sigma(A)$ are
$t$-degenerate for the order by position and its transport by $\sigma$ (\cref{lem:options,lem:rounding}); and the Lean
library proves the expected loads and second moments of \cref{lem:fluid} and the range bounds
\lean{AboveHalf.Fluid.f1_sub_le_one} and \lean{AboveHalf.Fluid.f0_sub_le_one}, which are the inputs of Bernstein's
inequality.

\begin{corollary}[The base for $p=127$]\label{cor:base}
\partformal{AboveHalfRate/Numbers.lean}{AboveHalf.Numbers.base_degeneracy,AboveHalf.Numbers.base_bernstein}{the application of \cref{thm:D}, whose probabilistic step is not formalized.}
Every simple $16130$-regular graph on $n=2\,197\,862$ vertices has a $4640$-degenerate lift selection.
\end{corollary}
\begin{outline}
Take $\delta=16130$ and $\lambda=317$ in \cref{thm:D}. Since $\tau(1-\tau)=3\sqrt3-5<0.1961527$, we have $\nu<3163.95$
and $\lambda^2/(2(\nu+\lambda/3))>100489/6539.24>15.367$, while $\log(2n)=\log4\,395\,724<15.297$. So \eqref{eq:Dcond}
holds, and $t=\lceil16130\tau+317\rceil=\lceil4639.02\ldots\rceil=4640$. The Lean theorems \lean{AboveHalf.Numbers.base_degeneracy} and \lean{AboveHalf.Numbers.base_bernstein} check
\eqref{eq:Dcond} and $t=4640\le4663$ exactly, through their own rational bounds on $\sqrt3$, on $e$ and on the
exponent; the decimals above are numerical evaluations.
\end{outline}

\begin{remark}[Why the order must be deep]\label{rem:naive}
\Cref{sec:theoremA} needs selections whose degeneracy is a fraction of $\delta$ below $1-1/\sqrt2\approx0.293$.
Orders with only two levels cannot give that. Suppose every lift in $A$ has an endpoint with at most $t$ lifts of $A$
(an \emph{$A$-light} vertex), every lift in $B$ an endpoint with at most $t$ lifts of $B$ (a \emph{$B$-light}
vertex), and $t<\delta/2$; then the private checks of \cref{sec:peeling} could all be single local checks. The swap
exchanges $A$-light and $B$-light vertices, so there are equally many, say $\ell$. A $B$-light vertex carries at least
$\delta-t>t$ lifts of $A$, and each has its other end $A$-light; an $A$-light vertex carries at most $t$. Hence
$\ell(\delta-t)\le\ell t$, so $\ell=0$, which is impossible because $A$ is not empty. So $t\ge\delta/2$, and the rate
$(1-t/\delta)^2$ of \cref{sec:theoremA} would be at most $1/4$. Deep orders, as in \cref{thm:D}, are what make a smaller
degeneracy possible, and they are why the private checks have large support. This remark is not used in any proof.
\end{remark}
